% Part: proof-theory
% Chapter: sequent-calculus
% Section: admissible-derivable

\documentclass[../../../include/open-logic-section]{subfiles}

\begin{document}

\olfileid[hi]{pt}{seq}{adm}

\olsection{ग्राह्य और व्युत्पाद्य नियम}

प्रमाण सिद्धांत के अनेक परिणाम इस प्रश्न से संबद्ध परिणामों से सरोकार रखते
हैं---या उनका उपयोग करते हैं---कि किसी विशिष्ट नियम का उपयोग करके !!{prove}
की जा सकने वाली हर वस्तु को क्या उस नियम के बिना भी !!{prove} किया जा सकता
है। कट-विलोपन प्रमेय इसका सबसे प्रसिद्ध उदाहरण है। वह दिखाता है कि~\Cut{}
सहित किसी अनुक्रम प्रणाली के नियमों से !!{prove} किए जा सकने वाले प्रत्येक
अनुक्रम को~\Cut के बिना भी !!{prove} किया जा सकता है। इस प्रकार का गुणधर्म
इतना महत्त्वपूर्ण है कि इसका अपना नाम है।

\begin{defn}
कोई नियम~$R$ किसी अनुक्रम प्रणाली में \emph{ग्राह्य} है यदि उस प्रणाली और
नियम~$R$ को साथ लेकर सिद्ध किए जा सकने वाले प्रत्येक अनुक्रम को~$R$ के बिना
भी सिद्ध किया जा सकता है।
\end{defn}

\begin{ex}\ollabel{ex:land-deriv}
नियम $\LeftR{\land}$ के \Log{G1c} और \Log{G3c} में अलग-अलग रूप हैं।
\Log{G1c} में इसके दो रूप हैं: एक में प्रमुख !!{formula}~$!A \land !B$ का
वाम संयोज्य $!A$ आधार-वाक्य में है, और दूसरे में दक्षिण संयोज्य~$!B$:
\[
\Axiom$ !A, \Gamma \fCenter \Delta$
\RightLabel{$\LeftR{\land}_1$}
\UnaryInf$ !A \land !B, \Gamma \fCenter \Delta$
\DisplayProof
\qquad
\Axiom$!B, \Gamma \fCenter \Delta$
\RightLabel{$\LeftR{\land}_2$}
\UnaryInf$!A \land !B, \Gamma \fCenter \Delta$
\DisplayProof
\]
इसके विपरीत, \Log{G3c} में केवल एक नियम है:
\[
\Axiom$ !A, !B, \Gamma \fCenter \Delta$
\RightLabel{$\LeftR{\land}_3$}
\UnaryInf$ !A \land !B, \Gamma \fCenter \Delta$
\DisplayProof
\]
अब हम पूछ सकते हैं, ``क्या \Log{G3c} का नियम $\LeftR{\land}_3$, \Log{G1c}
में ग्राह्य है?'' हाँ: $\LeftR{\land}_3$ का उपयोग करने वाले किसी भी निगमन का
हम केवल \Log{G1c} के नियमों से सीधे अनुकरण कर सकते हैं।
$\LeftR{\land}_1$ और $\LeftR{\land}_2$ के अतिरिक्त हमें संरचनात्मक
नियम~$\LeftR{\Contraction}$ भी चाहिए:
\[
\Axiom$!A, !B, \Gamma \fCenter \Delta$
\RightLabel{$\LeftR{\land}_1$}
\UnaryInf$!A \land !B, !B, \Gamma \fCenter \Delta$
\RightLabel{$\LeftR{\land}_2$}
\UnaryInf$!A \land !B, !A \land !B, \Gamma \fCenter \Delta$
\RightLabel{\LeftR{\Contraction}}
\UnaryInf$!A \land !B \Gamma \fCenter \Delta$
\DisplayProof
\]
\end{ex}

एक नियम वाले निगमन का अन्य नियमों से अनुकरण करना किसी नियम को ग्राह्य दिखाने
का एक ढंग है। किंतु यही एकमात्र ढंग नहीं है, और कुछ नियम ग्राह्य तो हैं पर
इस प्रकार अनुकरणीय नहीं। जिन नियमों का ऐसा अनुकरण हो सके उन्हें
\emph{व्युत्पाद्य} कहते हैं।

\begin{defn}
कोई नियम~$R$,
\[
\AxiomC{$S_1$}
\AxiomC{$\dots$}
\AxiomC{$S_n$}
\RightLabel{$R$}
\TrinaryInfC{$S$}
\DisplayProof
\]
किसी प्रणाली में \emph{व्युत्पाद्य} कहलाता है यदि प्रणाली के नियमों और
अभिगृहीतियों के साथ~$S_1$, \dots,~$S_n$ रूप के अतिरिक्त आरंभिक अनुक्रमों का
उपयोग करके अनुक्रम~$S$ का कोई योजनाबद्ध !!{proof} हो।
\end{defn}

\olref{ex:land-deriv} में दिया !!{proof} दिखाता है कि $\LeftR{\land}_3$,
\Log{G1c} में व्युत्पाद्य नियम है, क्योंकि वह केवल \Log{G1c} के नियमों से,
$\LeftR{\land}_3$ के आधार-वाक्यों को आरंभिक अनुक्रम मानकर, उसके निष्कर्ष का
योजनाबद्ध !!{proof} है। (उसमें \Log{G1c} की कोई अभिगृहीति उपयोग नहीं हुई।)

\begin{prop}\ollabel{prop:lor-G3-adm}
\Log{G3c} के नियम $\LeftR{\land}$ और $\RightR{\lor}$, \Log{G1c} में
व्युत्पाद्य हैं।
\end{prop}

\begin{proof}
  $\LeftR{\land}$ का प्रमाण \olref{ex:land-deriv} में दिया है।
  $\RightR{\lor}$ का प्रमाण अभ्यास के लिए छोड़ा गया है।
\end{proof}

\begin{prob}
दिखाएँ कि \Log{G3c} का नियम $\RightR{\lor}$, \Log{G1c} में व्युत्पाद्य है।
\end{prob}

\begin{prob}
मान लें कि $\liff$ कोई परिभाषित !!{operator} है: $!A \liff !B$, $(!A \lif
!B) \land (!B \lif !A)$ का संक्षिप्त रूप है। दिखाएँ कि नियम
\begin{enumerate}
\item \Axiom$\Gamma, !A, !B \fCenter \Delta$
\Axiom$\Gamma \fCenter \Delta, !A, !B$
\RightLabel{\LeftR{\liff}}
\BinaryInf$\Gamma \fCenter \Delta, !A \liff !B$
\DisplayProof
\item
\Axiom$!A, \Gamma \fCenter \Delta, !B$
\Axiom$!B, \Gamma \fCenter \Delta, !A$
\RightLabel{\RightR{\liff}}
\BinaryInf$\Gamma \fCenter \Delta, !A \liff !B$
\DisplayProof
\end{enumerate}
(क)~\Log{G1c} और (ख)~\Log{G3c} में व्युत्पाद्य हैं।
\end{prob}


\begin{ex}\ollabel{ex:weak-adm}
\Log{G1c} का नियम $\RightR{\Weakening}$,
\[
\Axiom$ \Gamma \fCenter \Delta$
\RightLabel{\RightR{\Weakening}}
\UnaryInf$ \Gamma \fCenter \Delta, !A$
\DisplayProof
\]
\Log{G3c} में ग्राह्य है। प्रमाण का विचार सरल है: मान लें कि \Log{G3c} में
आधार-वाक्य $\Gamma \Sequent \Delta$ का कोई !!a{proof}~$\pi$ है। $\pi$ के
हर अनुक्रम के अनुवर्ती में !!{formula}~$!A$ जोड़ दें। अभिगृहीतियाँ अभिगृहीतियाँ
रहती हैं; सही निगमन सही निगमन रहते हैं। नए !!{proof} का अंतिम-!!{formula}
$\Gamma \Sequent \Delta, !A$ है।

किंतु नियम $\RightR{\Weakening}$, \Log{G3c} में व्युत्पाद्य नहीं है। मान लें
कि वह है, अर्थात \Log{G3c} की केवल अभिगृहीतियों और नियमों तथा संभवतः अतिरिक्त
आरंभिक अनुक्रम~$\Gamma \Sequent \Delta$ से $\Gamma \Sequent \Delta, !A$ का
कोई !!a{proof} मिल सकता है। यदि यह !!{proof}, जैसा $\RightR{\Weakening}$ के
व्युत्पाद्य होने के लिए आवश्यक है, पूर्णतः योजनाबद्ध हो, तो $\Gamma$ और
~$\Delta$ हटाकर हम उसे~$\Sequent !A$ के !!a{proof} में बदल सकते हैं। इससे
अतिरिक्त आरंभिक अनुक्रम~$\Gamma \Sequent \Delta$, रिक्त अनुक्रम~$\quad
\Sequent \quad$ में बदल जाएगा। चूँकि रिक्त अनुक्रम में एक भी !!{formula} नहीं
है, जबकि \Log{G3c} के हर नियम को अपने आधार-वाक्यों में कम-से-कम एक सक्रिय
!!{formula} चाहिए, रिक्त अनुक्रम इस योजनाबद्ध !!{proof} के किसी निगमन का
आधार-वाक्य नहीं हो सकता। फलतः !!{proof} तभी योजनाबद्ध हो सकता है जब वह वास्तव
में अतिरिक्त अनुक्रम $\Gamma \Sequent \Delta$ का आरंभिक अनुक्रम के रूप में
उपयोग न करे। दूसरे शब्दों में, वह केवल \Log{G3c} की अभिगृहीतियों और नियमों से
$\Gamma \Sequent \Delta, !A$ की योजनाबद्ध व्युत्पत्ति होगा, और हमने दिखा दिया
होगा कि इस रूप के हर अनुक्रम का \Log{G3c} में !!a{proof} है। किंतु यह असंभव
है, क्योंकि \Log{G3c} निर्दोष है और इसलिए केवल उन अनुक्रमों को !!{prove} कर
सकता है जो हर !!{structure} में सत्य हैं। $\Gamma = \Delta = \emptyset$ और
$!A \ident \lfalse$ लेने पर \Log{G3c} को, उदाहरणतः, अनुक्रम $\quad \Sequent
\lfalse$ को !!{prove} करना होगा, जो वह नहीं करता।
\end{ex}

आइए इस परिणाम का एक सुंदर आगमनात्मक प्रमाण दें कि दुर्बलीकरण \Log{G3c} में
ग्राह्य नियम है। वास्तव में, जैसा सहज तर्क दिखाता है, थोड़ा अधिक प्रबल परिणाम
सत्य है: $\pi$ के हर अनुक्रम की दाईं ओर $!A$ जोड़ने से हम !!{proof} की संरचना
नहीं बदलते, विशेषतः उसकी ऊँचाई नहीं बढ़ाते। यह महत्त्वपूर्ण है, इसलिए इसकी
अपनी परिभाषा उचित है।

\begin{defn}
कोई नियम~$R$,
\[
\AxiomC{$S_1$}
\AxiomC{$\dots$}
\AxiomC{$S_n$}
\RightLabel{$R$}
\TrinaryInfC{$S$}
\DisplayProof
\]
किसी प्रणाली में \emph{!!{height}-संरक्षक ग्राह्य} (या
\emph{!!{hp}-ग्राह्य}) है यदि $i = 1$, \dots,~$n$ के लिए जब भी
$\Proves[h_i] S_i$ हो, तब $m = \max(h_1, \dots, h_n)$ के साथ $\Proves_m S$
हो।
\end{defn}

ध्यान दें कि किसी नियम के !!{height}-संरक्षक ग्राह्य होने के लिए इतना पर्याप्त
है कि जब भी आधार-वाक्यों~$S_i$ के !!{height}~$h_i = \pheight{\pi_i}$ वाले
!!{proof}s~$\pi_i$ हों, तब निष्कर्ष का कोई ऐसा !!a{proof}~$\pi$ हो जिसकी
!!{height}~$\pheight{\pi}$, $h_i$ के अधिकतम से बड़ी न हो। विशेषतः, यदि केवल
एक आधार-वाक्य~$S_1$ है, तो $\pheight{\pi} \le \pheight{\pi_1}$ पर्याप्त है।
$\RightR{\Weakening}$ और $\LeftR{\Weakening}$ की स्थिति में वास्तव में
$\pheight{\pi} = \pheight{\pi_1}$ है, किंतु यह आवश्यक नहीं।

\begin{prop}\ollabel{prop:weak-G3c-adm}
\Log{G1c} के नियम $\RightR{\Weakening}$ और $\LeftR{\Weakening}$,
\Log{G3c} में !!{height}-संरक्षक ग्राह्य हैं।
\end{prop}

\begin{proof}
  हमें दिखाना है कि यदि !!{height}~$\pheight{\pi} = n$ वाले !!{proof}~$\pi$
  के साथ $\Log{G3c} \Proves \Gamma \Sequent \Delta$, तो $\Gamma \Sequent
  \Delta, !A$ का कोई \Log{G3c}-!!{proof} $\pi'$ है जिसके लिए
  $\pheight{\pi'} = n$। हम इसे~$n$ पर आगमन द्वारा सिद्ध करते हैं। यदि $n=0$,
  तो $\pi$ केवल अभिगृहीति $\Gamma \Sequent \Delta$ है (अर्थात कोई परमाण्विक
  $!C$, $\Gamma$ और~$\Delta$ दोनों में आता है), और $\Gamma \Sequent \Delta,
  !A$ भी अभिगृहीति है। यदि $n = \pheight{\pi} >0$, तो !!{proof}~$\pi$ का कोई
  अंतिम निगमन है। उपयोग हुए नियम के अनुसार स्थितियाँ अलग करते हैं। केवल एक
  उदाहरण दें: मान लें कि अंतिम निगमन $\RightR{\land}$ था। तब $\Delta =
  \Delta', !B \land !C$, और अंतिम निगमन के निकटस्थ उप-!!{proof}s $\pi_1$ और
  ~$\pi_2$ के अंतिम अनुक्रम क्रमशः $\Gamma \Sequent \Delta', !B$ और $\Gamma
  \Sequent \Delta', !C$ हैं। दूसरे शब्दों में, $\pi$ ऐसा दिखता है:
  \[
  \AxiomC{}
  \RightLabel{$\pi_1$}
  \Deduce$\Gamma \fCenter \Delta', !B$
  \AxiomC{}
  \RightLabel{$\pi_2$}
  \Deduce$\Gamma \fCenter \Delta', !C$
  \RightLabel{\RightR{\land}}
  \BinaryInf$\Gamma \fCenter \Delta', !B \land !C$
  \DisplayProof
  \]
  हमारे पास $n = \max(\pheight{\pi_1}, \pheight{\pi_2}) + 1$ भी है। अतः
  $\pheight{\pi_1} < n$ और $\pheight{\pi_2} < n$, इसलिए आगमन परिकल्पना लागू
  होती है: $\Gamma \Sequent \Delta', !B, !A$ और $\Gamma \Sequent \Delta',
  !C, !A$ के !!{proof}s $\pi_1'$ और~$\pi_2'$ हैं। $\RightR{\land}$ नियम से
  हमें $\Gamma \Sequent \Delta', !B \land !C, !A$ का !!a{proof}~$\pi'$
  मिलता है:
  \[
  \AxiomC{}
  \RightLabel{$\pi_1'$}
  \Deduce$\Gamma \fCenter \Delta', !B, !A$
  \AxiomC{}
  \RightLabel{$\pi_2'$}
  \Deduce$\Gamma \fCenter \Delta', !C, !A$
  \RightLabel{\RightR{\land}}
  \BinaryInf$\Gamma \fCenter \Delta', !B \land !C, !A$
  \DisplayProof
  \]
  हमारे पास $\pheight{\pi'} =
  \max(\pheight{\pi_1'}, \pheight{\pi_2'}) + 1 = \max(\pheight{\pi_1},
  \pheight{\pi_2}) + 1 = \pheight{\pi}$.
\end{proof}

\olref{prop:weak-G3c-adm} अत्यंत उपयोगी है। इसके बार-बार अनुप्रयोग के लिए हम
एक संकेतन प्रस्तुत करेंगे: यदि $\pi$, $\Gamma \Sequent \Delta$ का
!!a{proof} है, तो $\pi[\Pi \Sequent \Lambda]$, बाईं ओर~$\Pi$ के सभी
!!{formula}s के साथ दुर्बलीकरण की ग्राह्यता और दाईं ओर~$\Lambda$ के सभी
!!{formula}s के साथ दुर्बलीकरण की ग्राह्यता लागू करने से मिला परिणाम है। यह
$\Gamma, \Pi \Sequent \Delta, \Lambda$ का !!a{proof} है।

\end{document}
